\section{Signs of Intelligence}

\begin{quotex}
\begin{verse}
Imagine, as philosophers postulate,\\
that the earth is another moon:\\
then we're all lunatic and vertiginous\\
within a sublunary maze. Can you doubt it?
\end{verse}
\flright{\textsc{Evan S Connell}, \emph{Points on a Compass Rose}}
\end{quotex}

\flrightit{Posted on 2012-02-01 by Cologero}

\begin{center}* * *\end{center}

\begin{footnotesize}
\begin{sffamily}
\texttt{Charlotte Cowell on 2012-02-02 at 05:01 said:}

excellent, best post yet, lol! and the answer is... .a GSOH :-))

\hfill

\texttt{HOO on 2012-02-02 at 07:30 said:}

Interestingly I find this to be one of the worst posts yet and not a sign of intelligence.

``As for life after death, the Vedic solution is fully consonant with the oldest Aryo-Hellenic spirit: images of obscure hells are almost entirely absent from the most ancient parts of the Vedas; the crisis of death is hardlynoticed as such-in the Atharva Veda it is even considered as the effect of a hostile and demoniacal force that, with suitable rites, can be repulsed.''

\hfill

\texttt{Cologero on 2012-02-02 at 08:53 said:}

And who, today, is performing those suitable rites?

\hfill

\texttt{uberman on 2012-02-02 at 09:51 said:}

Just as important a question is What are those rites?

A forth option for the last: the Individual who is still living yet has transcended death. The diamond thunder-bolt body of the immortal conqueror of the serpent hydra. Thus self-mastery complete.

\hfill

\texttt{HOO on 2012-02-02 at 09:53 said:}

Haha, I don't know, nor would we probably know if anyone was doing so as it seems a rather secluded matter. However I do intend to, and everything else on earth seems like a joke in comparison.

Anyway, ``There are, however, spiritual movements wherein people are tied by liberty. People are free to join this spiritual family. No exterior determination forces them to become brothers. Back in the day when it was expanding and converting, Christianity was a spiritual movement that people joined out of the common desire to spiritualize their lives and overcome death.'' --Mircea Eliade, 1937

``Nolite timere eos qui corpus occidunt, animum autem occidere non possunt'' Matt 10:22 ``Guénon interpreted the Latin word `Amor' as `a-mors' (without death)''

\hfill

\texttt{Gabe Ruth on 2012-02-02 at 10:42 said:}

I guess we're not intelligent, HOO.

\hfill

\texttt{Charlotte Cowell on 2012-02-02 at 11:37 said:}

For one reason or another I'm put in mind of an art history class from MANY moons ago, where we were studying ancient sculpture. Without organising an impromptu lecture I can recommend researching the meaning behind the `Archaic Smile'. I've even got a verse about it within the epic Myth:

\begin{verse}
Now does Hermes turn. Archaic\\
Is his smile. Osiris sees this,\\
Sees within it Thoth the Ancient\\
Science Master, Time Atomic.

Looks he, now -- the one -- at Hermes,\\
Thinks into his mind the verdict,\\
Weighs the words, the vital message.\\
Mount Olympus quaked to hear it.

I shall not deny the Certain\\
Things, these things shall be conceded --\\
Such as sharing life-eternal --\\
This, I think, by law, is needed.

Through the gift of sacred music,\\
Orphic guides shall spring forever\\
In the minds and hearts of humans,\\
Everflowing love around them.
\end{verse}

And so it goes on... ..

\url{http://alchemical-weddings.com/alchemical-weddings/archaic-smile}


\hfill

\texttt{Charlotte Cowell on 2012-02-02 at 11:50 said:}

Here's another one for you Hoo, seeing as you've migrated East (not a bad idea, I've often thought of moving to India, I love the place):

\begin{verse}
From the swelling seas, un-silent,\\
Rising from the salt, through ether,\\
Neptune holds aloft his trident,\\
Cries: ``The Spring has come; be patient!

As the centre of his offspring\\
Glows -- outraged to so be lectured --\\
So much wisdom of the ages\\
Flows from father-ocean's lectern:

``Take some good advice, Orion:\\
Watch and learn the way of heaven;\\
Time just moves around in circles,\\
From the fish becomes a turtle.''

``Onward then in time, a deluge\\
Caused a boar to swim the ocean;\\
Then the lion, Narasimha\\
Came before the dwarf Vamana.

``Then to life a noble hero\\
Sprang and rid the world of tyrants.\\
This made way for Rama's charm,\\
Which came before the Bhagavad Gita.

``In this way the prince of paupers\\
Broke the wheel of earthly suffering;\\
Maybe, son, you'll hear him teaching\\
In the realm of endless loving... ''
\end{verse}

\hfill

\texttt{Charlotte Cowell on 2012-02-02 at 11:55 said:}

Needless to say Orion refuses to listen and presses on with his plan to raise an army from the ranks of constellations, which does not succeed either... 

\begin{verse}
''Will you see, great Herakles\\
That we are stronger, still, together,\\
Strong enough to reach Elysium;\\
Gods are great, but we are better?
\end{verse}

However he DOES get his revenge sooner or later, but only when he wakes up, remembers who he is and is marvellously rescued by his... wife... ... 

\hfill

\texttt{uberman on 2012-02-02 at 19:57 said:}

HOO said: ``Haha, I don't know, nor would we probably know if anyone was doing so as it seems a rather secluded matter. However I do intend to, and everything else on earth seems like a joke in comparison.''

Yes it is really that simple, and I share your sentiment about the joke of it all. In essence, nothing matters but that which is directed towards the transcendent. All aims of life should point vertical from state to individual and all else besides is transitory and meaningless (at least it seems this way to me). In light of that understanding, why waste time in meddling amongst the yentas?

It brings more power to the thought that Evola had penned: (to paraphrase) ``Necessity is an evil''.

\hfill

\texttt{Charlotte Cowell on 2012-02-03 at 02:53 said:}

`meddling amongst the yentas'... ..

charming but typical :-)

\hfill

\texttt{HOO on 2012-02-11 at 11:18 said:}

``... the Devas, in their fight against the Asuras, protect themselves (achhandayan) by the recitation of the hymns of the Veda, and that it is for this reason that the hymns received the name of chhandas, a word which denotes ``rhythm.'' The same idea is contained in the word dhikr which, in Islamic esoterism, is used of rhythmic formulas that correspond exactly to Hindu mantras. The repetition of these formulas aims at producing a harmonization of the different elements of the being, and at causing vibrations which, by their repercussions throughout the immense hierarchy of states, are capable of opening up a communication with the higher states, which in a general way is the essential and primordial purpose of all rites.'' Guénon, ``The Language of the Birds'' \hfill

\texttt{Charlotte Cowell on 2012-02-11 at 13:25 said:}

Seeing as you mentioned birds I'm bound to draw attention to one of the most wonderful poems ever written, Attar's Bird Parliament. Here's a bit to remind us!

\begin{verse}
One night from out the swarming city gate\\
Stept holy Bajazyd, to meditate\\
Alone amid the breathing fields that lay\\
In solitary silence leagues away,\\
Beneath a Moon and Stars as bright as Day.\\
And the Saint wondering such a temple were,\\
And so lit up, and scarce one worshipper,\\
A voice from Heav'n amid the stillness said:\\
``The Royal Road is not for all to tread,\\
Nor is the Royal Palace for the rout,\\
Who, even if they reach it, are shut out.\\
The blaze that from my harim window breaks\\
With fright the rabble of the roadside takes;\\
And ev'n of those that at my Portal din,\\
Thousands may knock for one that enters in.
\end{verse}

\hfill

\texttt{HOO on 2012-02-11 at 18:45 said:}

Excellent poem.

\hfill

\texttt{White Rabbit on 2015-11-21 at 22:14 said:}

I imagine Dante in English gets ``lost in translation''

\hfill

\texttt{HOO on 2019-08-20 at 22:50 said:}

Aside from one false dichotomy, I must say that I have changed! The opening post is a good post.


\hfill

\end{sffamily}
\end{footnotesize}

